\section{Introduction}
\label{sec:introduction}

The intersection of physics-based modeling and machine learning has emerged as a frontier for scientific prediction, particularly in domains where first-principles understanding exists but computational or modeling limitations prevent accurate forecasting. Climate science presents a compelling testbed for such hybrid approaches: the underlying physics is well-established, yet dynamical models remain imperfect, and observational records are short relative to the timescales of interest. El Ni\~no-Southern Oscillation (ENSO), the dominant mode of interannual climate variability~\cite{mcphaden2006enso,timmermann2018enso}, exemplifies these challenges. Despite decades of research and operational forecasting efforts, predicting ENSO beyond six months remains difficult due to chaotic dynamics, model biases, and the so-called ``spring predictability barrier''~\cite{webster1995monsoon,duan2013predictability} where forecast skill drops sharply when predictions must traverse the boreal spring season.

Traditional approaches to ENSO forecasting span two paradigms. Dynamical models based on coupled ocean-atmosphere general circulation models (GCMs)~\cite{flato2013evaluation} encode comprehensive physical processes but suffer from systematic biases, require enormous computational resources, and struggle to outperform simpler statistical methods at lead times beyond one season~\cite{barnston2012skill}. Statistical models exploit observed relationships between climate variables but typically lack interpretability and may fail when climate statistics shift---a concerning limitation given anthropogenic climate change. Recent advances in deep learning have prompted efforts to apply neural networks directly to climate prediction~\cite{reichstein2019deep}, with notable successes in multi-year ENSO forecasting~\cite{ham2019deep}. However, these purely data-driven approaches face a fundamental obstacle: the observational record provides at most a few hundred independent samples for training, far fewer than the millions typically required for deep learning to excel.

This data scarcity motivates the development of hybrid physics-machine learning models that combine the interpretability and sample efficiency of physics-based approaches with the flexibility and representational power of neural networks~\cite{willard2022integrating,yin2021augmenting}. Recent work on Universal Differential Equations~\cite{rackauckas2020universal} has demonstrated that embedding neural networks within physics-based differential equation frameworks can discover missing dynamics while respecting known physical constraints. The central question we address is: \emph{How can we best integrate physical structure with learned components for climate forecasting when data is severely limited?}

We build upon the Extended Recharge Oscillator (XRO) model~\cite{zhao2024xro}, a recent physics-based framework that achieves state-of-the-art ENSO forecast skill by encoding the fundamental recharge-discharge dynamics~\cite{jin1997equatorial} of the tropical Pacific coupled with seasonal modulation and polynomial nonlinearities. While XRO provides strong baselines and interpretable parameters, its rigid polynomial basis may not capture all relevant dynamics, and its closed-form fitting procedure optimizes each variable independently rather than jointly for forecast skill.

Our contribution is the Neural Residual Ordinary Differential Equation (NXRO) framework, inspired by neural ODE methods~\cite{chen2018neuralode,rubanova2019latent} and physics-informed machine learning~\cite{karniadakis2021physics}. Unlike residual connections in deep networks~\cite{he2016deep} that provide skip connections between layers, our ``residual'' refers to learning the discrepancy between a physics-based model and observed data---a concept with roots in Bayesian model calibration~\cite{kennedy2001bayesian} and recently extended to differentiable physics~\cite{um2020solver,list2022learned}. This decomposition explicitly separates dynamics into a linear physics-based component and a learned nonlinear correction, providing strong inductive biases~\cite{bronstein2021geometric} that enable learning from limited data. We develop 43 model variants spanning a spectrum from pure physics (XRO) to pure machine learning (transformers), enabling systematic investigation of how much physical structure is necessary and where neural flexibility provides benefit.

Our key findings challenge conventional wisdom in several ways. First, we demonstrate that simpler hybrid architectures outperform both the full physics model and more complex neural variants. The best-performing model (NXRO-Res) combines only a seasonal linear operator with a residual multilayer perceptron, omitting the polynomial nonlinearities that are central to the physics-based XRO. Second, we show that joint gradient-based optimization of model parameters yields better out-of-sample generalization than closed-form fitting, even when the model structure is identical. Third, we provide a cautionary result for practitioners: pure transformer models, despite their success in natural language processing and computer vision, rank near the bottom of our evaluation, performing worse than the physics baseline. This contrasts with the success of transformers in weather forecasting~\cite{bi2023accurate,lam2023learning,pathak2022fourcastnet}, where training data spans decades of global observations at high temporal resolution. Our failure underscores the importance of physics-informed inductive biases when data is limited to a few hundred samples.

For probabilistic forecasting, we introduce a two-stage stochastic training procedure. Rather than fitting noise parameters post-hoc from model residuals or using external climate model ensembles, we optimize AR(1) noise parameters via maximum likelihood while keeping the dynamics model fixed. This approach achieves 1--13\% improvement in Continuous Ranked Probability Score (CRPS)~\cite{hersbach2000,gneiting2007} over XRO and produces better-calibrated uncertainty estimates.

The remainder of this paper proceeds as follows. Section~\ref{sec:background} provides background on ENSO dynamics and the XRO model formulation for readers unfamiliar with climate science. Section~\ref{sec:neural_ode} presents our NXRO architectures and training procedures. Section~\ref{sec:deterministic} describes our deterministic forecasting experiments with rigorous out-of-sample evaluation. Section~\ref{sec:stochastic} covers probabilistic forecasting and the two-stage stochastic training approach. Section~\ref{sec:conclusion} summarizes our findings, and Section~\ref{sec:future} discusses limitations and future directions.
